\section{Постановка задачи}
\label{sec:Chapter1} \index{Chapter1}

\subsection{Модель искажения сигнала}

Пусть $x(t)$ --- чистый речевой сигнал студийного качества, дискретизированный
с частотой $f_s$. В реальных условиях наблюдается искажённый сигнал $y(t)$,
который получается из $x(t)$ в результате последовательного действия нескольких
деградаций. Обобщённая модель искажения записывается как
\begin{equation}
    y(t) = \mathcal{C}\Big( \big(x * h\big)(t) + n(t) \Big),
    \label{eq:degradation}
\end{equation}
где $h(t)$ --- импульсная характеристика помещения (RIR), описывающая
реверберацию; $*$ --- операция свёртки; $n(t)$ --- аддитивный шум; а
$\mathcal{C}(\cdot)$ --- нелинейный оператор, моделирующий искажения
микрофонного тракта и сжатие аудиокодеком. Уровень аддитивного шума
характеризуется отношением сигнал/шум
\begin{equation}
    \mathrm{SNR} = 10 \lg \frac{\sum_t \big((x*h)(t)\big)^2}{\sum_t n^2(t)}
    \quad [\text{дБ}].
    \label{eq:snr}
\end{equation}

\subsection{Формальная постановка}

Требуется построить отображение $\mathcal{F}_\theta$ с параметрами $\theta$,
которое по искажённому сигналу $y$ восстанавливает оценку чистого сигнала
\begin{equation}
    \hat{x} = \mathcal{F}_\theta(y),
    \label{eq:restore}
\end{equation}
максимально близкую к $x$ по перцептивному качеству. Прямая минимизация ошибки
во временной области неинформативна, поскольку перцептивно близкие сигналы могут
сильно различаться пофазно. Поэтому, следуя парадигме
Miipher-2~\cite{miipher2}, задача переносится в пространство признаков
самоконтролируемого энкодера $E$. Обозначим $z = E(y)$ --- последовательность
признаков искажённого сигнала, $z^{\ast} = E(x)$ --- признаки чистого сигнала.
Система состоит из трёх компонентов (рис.~\ref{fig:pipeline}):
\begin{enumerate}
    \item \textbf{SSL-энкодер} $E$ --- замороженная предобученная модель,
        отображающая волновую форму в последовательность скрытых представлений;
    \item \textbf{очиститель признаков} $G_\phi$ --- обучаемый модуль,
        предсказывающий чистые признаки $\hat{z} = G_\phi(z) \approx z^{\ast}$;
    \item \textbf{нейросетевой вокодер} $V_\psi$ --- модуль, синтезирующий
        волновую форму $\hat{x} = V_\psi(\hat{z})$ из очищенных признаков.
\end{enumerate}

Очиститель обучается минимизировать рассогласование признаков
\begin{equation}
    \mathcal{L}_{\text{fc}}(\phi) =
    \big\| G_\phi(E(y)) - E(x) \big\|_1 +
    \lambda \big\| G_\phi(E(y)) - E(x) \big\|_2^2,
    \label{eq:fcloss}
\end{equation}
а вокодер --- восстанавливать волновую форму по признакам чистого сигнала.

\begin{figure}[H]
    \centering
    \begin{tikzpicture}[
        node distance=0.9cm and 1.0cm,
        block/.style={draw, rounded corners, minimum height=1.0cm,
            minimum width=2.2cm, align=center, font=\small},
        frozen/.style={block, fill=black!6},
        trained/.style={block, fill=black!16},
        >={Stealth[length=2.5mm]}]
        \node (in) [align=center, font=\small] {Искажённый\\сигнал $y$};
        \node (enc) [frozen, right=of in] {SSL-энкодер\\$E$ (заморожен)};
        \node (fc)  [trained, right=of enc] {Очиститель\\признаков $G_\phi$};
        \node (voc) [trained, right=of fc] {Вокодер\\$V_\psi$};
        \node (out) [align=center, font=\small, right=of voc] {Чистый\\сигнал $\hat{x}$};
        \draw[->] (in) -- (enc);
        \draw[->] (enc) -- node[above, font=\scriptsize] {$z$} (fc);
        \draw[->] (fc) -- node[above, font=\scriptsize] {$\hat{z}$} (voc);
        \draw[->] (voc) -- (out);
    \end{tikzpicture}
    \caption{Обобщённая схема системы восстановления речи в пространстве
        SSL-признаков}
    \label{fig:pipeline}
\end{figure}

\subsection{Техническое задание и требования к решению}

На основе сформулированной модели определены следующие требования к
разрабатываемой системе.

\paragraph{Функциональные требования.}
\begin{enumerate}
    \item Система должна принимать на вход монофонический речевой сигнал с
        частотой дискретизации не ниже 16~кГц и возвращать восстановленный
        сигнал той же длительности.
    \item Система должна одновременно подавлять аддитивный шум, реверберацию и
        артефакты кодеков, описанные моделью~\eqref{eq:degradation}.
    \item Все компоненты системы должны быть основаны исключительно на открытых
        предобученных моделях и наборах данных, допускающих воспроизведение.
    \item Система не должна требовать текстовой расшифровки или информации о
        дикторе на этапе вывода (в отличие от Miipher-1).
\end{enumerate}

\paragraph{Требования к качеству.} Восстановленный сигнал $\hat{x}$ должен
демонстрировать статистически значимое улучшение относительно искажённого входа
по объективным метрикам: PESQ~\cite{pesq}, STOI~\cite{stoi},
SI-SDR~\cite{sisdr} и неинтрузивной оценке DNSMOS~\cite{dnsmos}. Целевые
ориентиры: прирост PESQ не менее $0{,}6$ балла и STOI не менее $0{,}05$
относительно входного сигнала на тестовом наборе.

\paragraph{Требования к вычислительной эффективности.} Система должна работать
быстрее реального времени на одном GPU, то есть коэффициент реального времени
$\mathrm{RTF} = T_{\text{обр}} / T_{\text{ауд}} < 1$, где $T_{\text{обр}}$ ---
время обработки, $T_{\text{ауд}}$ --- длительность аудио. Суммарное число
обучаемых параметров (очиститель и вокодер) не должно превышать характерных для
открытых вокодеров значений ($\sim$15~млн).

\paragraph{Критерий успешности.} Задача считается решённой, если среди
исследованных конфигураций найдена хотя бы одна, удовлетворяющая всем
требованиям, и определена конфигурация, обеспечивающая наилучший компромисс
между качеством и быстродействием.
